\chapter{Putting It Together: A Queueing Model of Disaggregated Serving}\label{sec:L5}

The first four lectures built tools. This lecture assembles them into a model of a
disaggregated deployment and derives what a scheduler for it needs: what a miss costs
at each place where it can happen, what each instance should keep in memory and for how
long, how many sessions to admit, and when misses feed
themselves. Every step names the result of Lectures 1--4 that carries it.

\section{The model on one page}\label{L5:sec:model}

\Cref{fig:deployment} in \cref{L1:sec:model} shows the deployment, and
\cref{L1:def:workload,L1:def:costs} and the program of \cref{L1:ex:program} define the system exactly. We
recall the notation.

\begin{definition}[Disaggregated replica]\label{L5:def:replica}
A \term{disaggregated replica} consists of the following stations, visited in this
order by every turn.
\begin{enumerate}[label=(\alph*)]
\item A \emph{prefill instance} with a memory pool of $M_P$ bytes: a waiting turn is
  admitted, in arrival order, only when its whole context fits; the part of the pool not
  held by admitted turns keeps a \emph{prefix cache} of paused sessions, evicted least
  recently used first; admitted turns are prefilled one at a time, in arrival order, by
  a server of full capacity.
\item A \emph{link} that carries the turn's KV to a decode instance, shared equally by
  concurrent transfers: a processor-sharing station of capacity 1 (\cref{def:ps}).
\item A \emph{decode instance} with a pool of $M_D$ bytes, the same whole-context
  admission, and a processor-sharing batch of capacity $\varphi(m)$ for $m$ running turns.
\item The \emph{tool call}, a delay station (\cref{L4:def:delay}) with mean think time
  $Z$, after which the paused session sends its next turn with probability $p$.
\end{enumerate}
\end{definition}

\paragraph{Costs.} Prefilling $n$ new tokens on a cached prefix of $K$ tokens takes
\begin{equation}\label{L5:eq:prefill}
  P(n,K)=a\,n+b\,n\,(K+n/2)
\end{equation}
seconds. A \term{hit} (the session's prefix is cached) costs $S^{\mathrm{hit}}=P(n,K)$; a
\term{miss} costs $S^{\mathrm{miss}}=P(K+n,0)$, so the extra work of a miss is
$\Delta S=aK+bK^2/2$ (\cref{cor:price-mono}). We write $S$ for the prefill work of a turn,
$h$ for the hit rate and $\rho_P=\lambda\E[S]$ for the prefill load, with $\lambda$ the turn
rate. A transfer of $T_i$ tokens needs
\begin{equation}\label{L5:eq:transfer}
  X_i=x_0+\frac{\kappa\,T_i}{B}
\end{equation}
seconds of the link, with $x_0$ a fixed overhead, $\kappa$ the KV bytes per token and $B$
the bandwidth; the link load is $\rho_L=\lambda\E[X]$. Decoding $o_i$ tokens is a demand
$D_i$ at the decode batch. If every context needs about $c$ bytes, instance
$j\in\{P,D\}$ has $m_j=\lfloor M_j/c\rfloor$ memory slots, and its admission station is the
multi-server queue of \cref{prop:mmm} with service time equal to the holding time
(prefill plus transfer at the prefill instance, decode at the decode instance).

\paragraph{What the time to first token contains.} The first token is produced by the
prefill instance, so
\[
  \mathrm{TTFT}=W^M_P+W_P+S,
\]
the wait for a memory slot at the prefill instance, the wait for the prefill server,
and the prefill itself. The transfer and the decode do not enter it: they delay the
second token and every later one. They matter for the TTFT only through memory, since a
slot at the prefill instance stays held until the transfer has finished.

\section{Where the variance of prefill work comes from}\label{L5:sec:variance}

The Pollaczek--Khinchine formula makes the second moment of $S$ a first-class quantity
(\cref{cor:variance}). The law of total variance splits it into a part the scheduler
does not control and a part it does.

\begin{lemma}[Variance split of the hit/miss mixture]\label{L5:lem:split}
Let $S_h$ be the mixture of \cref{L3:def:mixture}: $S^{\mathrm{hit}}$ with
probability $h$ and $S^{\mathrm{miss}}$ with probability $1-h$, with means
$\mu_{\mathrm{hit}},\mu_{\mathrm{miss}}$ and variances
$\sigma^2_{\mathrm{hit}},\sigma^2_{\mathrm{miss}}$. Then
\[
  \Var[S_h]=\underbrace{h\,\sigma^2_{\mathrm{hit}}+(1-h)\,\sigma^2_{\mathrm{miss}}}_{\text{spread of the appends}}
  +\underbrace{h(1-h)\,(\mu_{\mathrm{miss}}-\mu_{\mathrm{hit}})^2}_{\text{hit/miss mixture}} .
\]
\end{lemma}

\begin{proof}
Let $I$ be the indicator of a hit. By the law of total variance,
$\Var[S_h]=\E[\Var(S_h\mid I)]+\Var(\E[S_h\mid I])$. The first term is
$h\sigma^2_{\mathrm{hit}}+(1-h)\sigma^2_{\mathrm{miss}}$. The conditional
mean $\E[S_h\mid I]$ takes the value $\mu_{\mathrm{hit}}$ with probability
$h$ and $\mu_{\mathrm{miss}}$ with probability $1-h$, a two-point
variable whose variance is $h(1-h)(\mu_{\mathrm{miss}}-\mu_{\mathrm{hit}})^2$.
\end{proof}

The first term comes from the spread of tool outputs (a line or a whole file) and is
not the scheduler's to change. The second is the KV policy's lever: it vanishes at
$h\in\{0,1\}$ and grows with the squared gap between miss and hit work, which by
\eqref{L5:eq:prefill} grows like $K^4$ once attention dominates. What a scheduler
must price is therefore not a ratio of variances but the rise of
$\lambda\E[S^2]/(2(1-\rho_P))$ that one more miss causes.

\section{Three prices}\label{L5:sec:prices}

A decision that changes where a paused session's KV lives changes the work of its next
turn at one of three stations.

\paragraph{(a) A miss at the prefill instance.} The prefill instance is an M/G/1 FIFO
queue with the whole server available to prefill, so \cref{thm:price} applies as it
stands: turning a fraction $q_i$ of arrivals from hits into misses raises the mean number
at the prefill instance by an amount in
$[\lambda\sum_iq_i\Phi_i,\ \frac{1-\rho_P}{1-\rho_P'}\lambda\sum_iq_i\Phi_i]$, with
\[
  \Phi_i=\Delta S_i
   +\frac{\lambda\bigl((S_i^{\mathrm{miss}})^2-(S_i^{\mathrm{hit}})^2\bigr)}{2(1-\rho_P)}
   +\frac{\lambda W_P\,\Delta S_i}{1-\rho_P}.
\]
The middle term is the head-of-line term: every turn that arrives while the miss is
prefilled waits for all of it.

\paragraph{(b) A miss at the decode instance.} If the decode instance still holds the
session's KV from its previous turn, the next turn transfers only its new tokens,
$T_i=n_i$; if not, it transfers the whole context, $T_i=K_i+n_i$. A decode-side miss
therefore adds $\Delta X_i=\kappa K_i/B$ of link work.

\begin{proposition}[The price of a transfer miss]\label{L5:prop:link-price}
Let the link be a processor-sharing station of capacity 1 with load $\rho_L<1$, and
turn a fraction $q_i$ of arrivals into decode-side misses, raising the load to
$\rho_L'=\rho_L+\lambda\sum_iq_i\Delta X_i<1$. Then the mean number of transfers in
progress rises exactly by
\[
  \Delta L_L=\frac{1-\rho_L}{1-\rho_L'}\;\lambda\sum_iq_i\,\Phi^X_i,\qquad
  \Phi^X_i=\frac{\Delta X_i}{(1-\rho_L)^2},
\]
whatever the distribution of the transfer times.
\end{proposition}

\begin{proof}
By insensitivity (\cref{thm:insens}) the number of transfers in progress has the
geometric law of \cref{prop:decode-price} with capacity $C=1$, for every transfer-time
distribution with the given mean. \Cref{prop:decode-price} with $C=1$ and $\Delta D_i$
replaced by $\Delta X_i$ gives the formula.
\end{proof}

\paragraph{(c) Decode demand.} A decision that adds decode demand, such as admitting a
session, pays the load factor of \cref{prop:decode-price} at the decode batch; a miss at
either instance adds none (\cref{L2:cor:decode-hit}).

\paragraph{Comparison.} The same context of $K$ tokens costs the prefill instance a
miss whose work is quadratic in $K$ and whose price has a head-of-line term, set by the
second moment of the prefill work; it costs the link a transfer linear in $K$ whose
price depends on its mean only, because under processor sharing a long transfer does not
block the short ones. The prefix cache of the prefill instance is where memory buys the
most.

\section{The decision problem with two pools}\label{L5:sec:decision}

At a decision epoch each instance holds KV for a set of paused sessions. Session $i$
holds $c_i$ bytes, resumes with probability $p_i$ after an expected further $\tau_i$
seconds, and its next turn pays $\Phi_i$ if its prefix is not in the prefill instance's
cache and $\Phi^X_i$ if its KV is not at the decode instance. With indicators
$x^P_i,x^D_i\in\{0,1\}$ of ``keep'' at each instance and byte-second budgets $B_P$ and
$B_D$, the keep-or-drop problem is
\begin{equation}\label{L5:eq:primal}
  \mathrm{OPT}=\min_{x^P,x^D}\ \sum_i\bigl[(1-x^P_i)\,p_i\Phi_i+(1-x^D_i)\,p_i\Phi^X_i\bigr]
  \quad\text{s.t.}\quad \sum_ic_i\tau_ix^P_i\le B_P,\ \ \sum_ic_i\tau_ix^D_i\le B_D .
\end{equation}
Relaxing the two constraints with multipliers $\theta_P,\theta_D\ge0$, the
\term{shadow prices} of a byte-second at each instance, gives the Lagrangian
$\mathcal{L}=\sum_i[(1-x^P_i)p_i\Phi_i+(1-x^D_i)p_i\Phi^X_i]
+\theta_P(\sum_ic_i\tau_ix^P_i-B_P)+\theta_D(\sum_ic_i\tau_ix^D_i-B_D)$.

\begin{proposition}[One price per pool]\label{L5:prop:shadow}
Let $u^P_i=p_i\Phi_i/(c_i\tau_i)$ and $u^X_i=p_i\Phi^X_i/(c_i\tau_i)$.
\begin{enumerate}[label=(\roman*)]
\item For all $\theta_P,\theta_D\ge0$, $\mathcal{L}$ is minimised by keeping session $i$ at
  the prefill instance iff $u^P_i>\theta_P$ and at the decode instance iff
  $u^X_i>\theta_D$, and the minimum $g(\theta_P,\theta_D)$ is at most $\mathrm{OPT}$.
\item If the multipliers are such that these choices use both budgets exactly, the
  choices are optimal for \eqref{L5:eq:primal}.
\end{enumerate}
\end{proposition}

\begin{proof}
(i) $\mathcal{L}=\sum_ip_i(\Phi_i+\Phi^X_i)-\theta_PB_P-\theta_DB_D
+\sum_ix^P_i(\theta_Pc_i\tau_i-p_i\Phi_i)+\sum_ix^D_i(\theta_Dc_i\tau_i-p_i\Phi^X_i)$ is
separable: each $x^P_i$ is best set to 1 exactly when $\theta_Pc_i\tau_i<p_i\Phi_i$, that
is $u^P_i>\theta_P$, and likewise each $x^D_i$. For a feasible $(x^P,x^D)$ both brackets
multiplying $\theta_P$ and $\theta_D$ are nonpositive, so its objective is at least its
Lagrangian, which is at least $g$; minimising over feasible choices gives
$\mathrm{OPT}\ge g$ (weak duality). (ii) The choices of (i) are then feasible and their
constraint terms vanish, so their objective equals $g\le\mathrm{OPT}$.
\end{proof}

This is the threshold rule of \cref{thm:threshold}, once per pool: each instance keeps
what is worth more per byte-second than its own memory.

\begin{corollary}[Free transfer]\label{L5:cor:free-transfer}
As $B\to\infty$, $\Phi^X_i\to0$ for every session, so for any $\theta_D>0$ the decode
instance keeps nothing for paused sessions, and all of $M_D$ serves running decodes.
\end{corollary}

\begin{proof}
$\Delta X_i=\kappa K_i/B\to0$, and $\rho_L\to\lambda x_0$ stays below 1 when $\lambda x_0<1$,
so $\Phi^X_i=\Delta X_i/(1-\rho_L)^2\to0$ and $u^X_i\to0<\theta_D$.
\end{proof}

A decode instance's cache is worth exactly the transfers it saves; the prefill
instance's cache is worth the recomputation it saves and the queueing that recomputation
would cause.

\section{How long to keep}\label{L5:sec:hold}

\Cref{L5:prop:shadow} decides at one instant. Because the Lagrangian is separable, once
the multiplier of an instance is fixed each paused session can also be treated on its own
over time, at each instance separately, with that instance's price $\Phi$ (prefill:
$\Phi_i$; decode: $\Phi^X_i$) and multiplier $\theta$ ($\theta_P$ or $\theta_D$): keep its state until some time $T$ and drop it then if the
program has not resumed. Let the program resume with probability $p\in(0,1]$, and, given
that it resumes, after a time with a continuous distribution function $G$ and density
$g$ on $[0,\infty)$. Holding the state costs $\theta c$ per second; dropping it costs
$\Phi$ if the program resumes later. The expected cost of the rule ``hold until $T$'' is
\begin{equation}\label{L5:eq:hold-cost}
  C(T)=\theta c\int_0^T\bigl(1-pG(t)\bigr)\dd t+\Phi\,p\,\bigl(1-G(T)\bigr),
\end{equation}
since the state is held at time $t$ exactly when the program has not resumed by $t$,
and a miss occurs exactly when it resumes after $T$. The quantity that decides is the
\term{resume hazard}
\[
  \eta(t)=\frac{p\,g(t)}{1-pG(t)},
\]
the rate at which a program that has not resumed by $t$ resumes at $t$.

\begin{proposition}[Optimal holding time]\label{L5:prop:hold}
Let $r=\theta c/\Phi$ with $\Phi>0$, and suppose $g$ is continuous.
\begin{enumerate}[label=(\roman*)]
\item $C'(T)=\bigl(1-pG(T)\bigr)\bigl(\theta c-\Phi\,\eta(T)\bigr)$ for every $T\ge0$.
  Holding a little longer lowers the cost exactly when $\Phi\,\eta(T)>\theta c$.
\item If $\eta$ is nonincreasing, let $T^*=\inf\{T\ge0:\eta(T)\le r\}$. If
  $T^*<\infty$ it minimises $C$ over $[0,\infty)$; if $T^*=\infty$, $C$ is nonincreasing
  and the state should be held until the program resumes. In particular $T^*=0$, drop at
  once, if and only if $\eta(0)\le r$.
\item If $G$ is exponential with mean $\tau$, then $\eta$ is nonincreasing and
  $\eta(0)=p/\tau$, so the state should be kept at all if and only if
  \[
    u=\frac{p\,\Phi}{c\,\tau}>\theta ,
  \]
  the price per byte-second of \cref{L5:prop:shadow}. If $p=1$ and $u>\theta$, the
  state is held until the program resumes. If $p<1$ and $u>\theta$, the optimal holding
  time is
  \[
    T^*=\tau\,\ln\Bigl[\frac{p}{1-p}\Bigl(\frac{\Phi}{\theta c\tau}-1\Bigr)\Bigr].
  \]
\end{enumerate}
\end{proposition}

\begin{proof}
(i) The integrand in \eqref{L5:eq:hold-cost} is continuous, so by the fundamental
theorem of calculus the first term has derivative $\theta c(1-pG(T))$, and the second
has derivative $-\Phi p\,g(T)$. Factoring out $1-pG(T)$ gives the formula when
$pG(T)<1$; when $pG(T)=1$ (possible only if $p=1$) both sides vanish, because
$g(T)=0$ wherever $G$ has reached $1$ and $g$ is continuous. The factor $1-pG(T)$ is
nonnegative, so the sign of $C'(T)$ is the sign of $\theta c-\Phi\eta(T)$.

(ii) For $T<T^*$, $\eta(T)>r$, so $C'(T)\le0$ by (i). For $T>T^*$, $\eta(T)\le r$
because $\eta$ is nonincreasing, so $C'(T)\ge0$. Hence $C$ is nonincreasing on
$[0,T^*]$ and nondecreasing on $[T^*,\infty)$, and $T^*$ is a minimiser. If
$T^*=\infty$, then $C'\le0$ everywhere. The last claim is the definition of $T^*$ at
$T=0$, using that $\eta$ is nonincreasing.

(iii) With $x=e^{-T/\tau}$, $g(T)=x/\tau$ and $1-pG(T)=1-p+px$, so
\[
  \eta(T)=\frac{p}{\tau}\cdot\frac{x}{1-p+px}.
\]
The map $x\mapsto x/(1-p+px)$ is nondecreasing on $(0,1]$, and $x$ decreases in $T$,
so $\eta$ is nonincreasing; it is the constant $1/\tau$ when $p=1$. At $T=0$,
$\eta(0)=p/\tau$, and $\eta(0)>r$ is $p\Phi/(c\tau)>\theta$. If $p=1$, $\eta\equiv1/\tau>r$,
so $T^*=\infty$ by (ii). If $p<1$, $\eta(T)\to0$ as $T\to\infty$, so $T^*$ is finite and
solves $\eta(T^*)=r$ by continuity:
$px/\tau=r(1-p+px)$, that is $x=r(1-p)/\bigl(p(1/\tau-r)\bigr)$, and
$T^*=-\tau\ln x=\tau\ln\bigl[\frac{p}{1-p}\bigl(\frac{1}{r\tau}-1\bigr)\bigr]$ with
$1/(r\tau)=\Phi/(\theta c\tau)$.
\end{proof}

So the instantaneous test $u_i>\theta$ of \cref{L5:prop:shadow} is exactly the optimal
decision to keep a state for exponential resume times, and part~(iii) adds how long to keep it.
For example, with $p=0.5$, $\tau=10$\,s, $\Phi=20$\,s, $c=1$ and $\theta=0.3$, $u=1>\theta$,
the optimal holding time is $T^*\approx17.3$\,s, and it lowers the expected cost from
$C(0)=p\Phi=10$ to $C(T^*)\approx5.6$. When the gap distribution has a decreasing hazard, as
heavy-tailed gaps typically do, part~(ii) gives a program-specific time-to-live.

\begin{remark}[Unknown resume times and the oracle]\label{L5:rem:ski}
If nothing is known about $G$, holding the state is the classical \emph{ski rental}
problem: pay rent $\theta c$ per second or buy once at price $\Phi$
\cite{karlin1988}. The rule ``hold for $\Phi/(\theta c)$ seconds'' costs at most twice
the cost in hindsight for any resume time: if the program resumes at
$\tau\le\Phi/(\theta c)$ it pays $\theta c\tau$, which is optimal; otherwise it pays
$2\Phi$ against a hindsight optimum of $\Phi$. The hindsight optimum, which knows
whether and when each program resumes, is called the \term{oracle}; it keeps a state
exactly when the program resumes and $\theta c\tau<\Phi$. It cannot be run online, but
it is the benchmark against which a policy's cost is reported as cost/OPT. The factor
two needs the program to resume: against a program that never does, the oracle pays
nothing, and no rule that holds the state for a positive time has a bounded ratio.
\end{remark}

\Cref{L5:prop:hold} holds $\theta$ and $\Phi$ fixed. In the real system both move with
the policy: what is dropped changes the load, and the load changes every $\Phi_i$ and
the value of $\theta$ that fills the pool. The joint optimum is a Markov decision problem
over the whole state, with no closed form; \cref{L5:prop:hold} is its solution at a
fixed point of $(\theta,\Phi)$.

\section{How many sessions to admit}\label{L5:sec:admit}

Two memory quantities decide admission. The running turns at each instance occupy slots:
with exponential holding times the wait for a slot is the Erlang C wait of
\cref{prop:mmm}, and it grows without bound as the offered load of slots approaches
$m_j$. The paused sessions occupy the prefix cache: with sessions arriving as a Poisson
process of rate $\Lambda$, living for $T_s$ and holding $K(t)$ tokens at age $t$, the mean
resident KV is $\kappa\Lambda\E[\int_0^{T_s}K(t)\dd t]$ (\cref{cor:resident-kv}), which
exceeds the naive product whenever long contexts belong to long sessions. The admission
step caps the live sessions near the number at which this estimate fills the prefix
cache, and holds the rest in an entry queue. The cap has a price of its own, the wait at
the entry queue: a session is worth admitting while the memory it will hold, priced at
$\theta_P$ per byte-second, costs less than that wait.

\section{Misses that feed themselves}\label{L5:sec:feedback}

So far the hit rate was an input. It is also an output. A miss re-inserts a whole
context into the prefill instance's pool, which evicts other paused prefixes sooner
(\emph{pool channel}); and a turn that waits longer is absent longer, which exposes its
own prefix to eviction (\emph{wait channel}). Let $W(h)$ be the mean wait at the prefill
instance at hit rate $h$, nonincreasing in $h$ (\cref{cor:mixture}), and let $G(T;h)$ be
the probability that a paused prefix survives an absence $T$ when the instance runs at
hit rate $h$, nonincreasing in $T$ and nondecreasing in $h$.

\begin{definition}[Feedback map]\label{L5:def:feedback}
The \term{feedback map} is $H(h)=\E_Z[\,G(Z+W(h);h)\,]$ on $[0,1]$, and an
\term{equilibrium} is a hit rate with $H(h)=h$.
\end{definition}

\paragraph{The wait channel is an admission rule.} Whether a waiting turn's prefix
can be evicted is decided by the engine, not by the workload. An engine that
takes the prefix out of the cache only when it schedules the turn (vLLM admits a
waiting request at the start of an iteration with the budget left, and until then
its cached blocks stay evictable) has the wait channel. An engine that pins the
prefix when the turn arrives has survival $G(Z;h)$ instead of $G(Z+W(h);h)$, a larger
value for every $h$, but it holds the pinned bytes and so lowers $G$ for the other
paused prefixes; the comparison of \cref{L5:prop:tarski}(ii) applies only when the
first effect dominates. In a simulation of the vLLM engine replaying the
short-context agent trace on the A100 pool of 128k tokens, pinning moved the cliff:
at one session every 3\,s the vLLM rule collapsed (full-hit rate 0.24) and the pinning
rule did not (0.80); at 3.5\,s and above the two differed by less than two points.

$H$ is nondecreasing: if $h\le h'$ then $W(h)\ge W(h')$, so
$G(Z+W(h);h)\le G(Z+W(h');h)\le G(Z+W(h');h')$ for every $Z$, and taking expectations
preserves the order. No continuity is needed for what follows.

\begin{proposition}[Extremal equilibria]\label{L5:prop:tarski}
Let $H:[0,1]\to[0,1]$ be nondecreasing.
\begin{enumerate}[label=(\roman*)]
\item $\bar h=\sup\{x:x\le H(x)\}$ is an equilibrium, and every equilibrium is at most
  $\bar h$; $\underline h=\inf\{x:H(x)\le x\}$ is an equilibrium, and every
  equilibrium is at least $\underline h$.
\item If $K:[0,1]\to[0,1]$ is nondecreasing and $H\le K$ pointwise, then
  $\bar h_H\le\bar h_K$ and $\underline h_H\le\underline h_K$.
\end{enumerate}
\end{proposition}

\begin{proof}
(i) Let $A=\{x\in[0,1]:x\le H(x)\}$; $0\in A$, so $\bar h=\sup A$ exists. For $x\in A$,
$x\le H(x)\le H(\bar h)$ because $x\le\bar h$ and $H$ is nondecreasing; so $H(\bar h)$ is
an upper bound of $A$ and $\bar h\le H(\bar h)$. Then $H(\bar h)\le H(H(\bar h))$, so
$H(\bar h)\in A$ and $H(\bar h)\le\bar h$. Hence $H(\bar h)=\bar h$. Every equilibrium lies
in $A$, so it is at most $\bar h$. The argument for $\underline h$ is the mirror image
with $B=\{x:H(x)\le x\}\ni1$ and infima.
(ii) $\bar h_H=H(\bar h_H)\le K(\bar h_H)$, so $\bar h_H$ lies in the set $A$ of $K$, and
$\bar h_H\le\bar h_K$. Likewise $H(\underline h_K)\le K(\underline h_K)=\underline h_K$
puts $\underline h_K$ in the set $B$ of $H$, so $\underline h_H\le\underline h_K$.
\end{proof}

A larger prefix cache raises $G$ and so $H$, and by (ii) it raises both extremal
equilibria. Other changes need care in a closed loop: a shorter think time or a faster
decode also raises the turn rate, and the net effect on $H$ has no sign in general.

\begin{proposition}[Collapse]\label{L5:prop:collapse}
In the open prefill instance take the wait of \cref{thm:pk}, $W(h)=\infty$ where
$\lambda\E[S_h]\ge1$, and $G(\infty;h)=0$.
\begin{enumerate}[label=(\roman*)]
\item If $\lambda\,\E[S^{\mathrm{miss}}]\ge1$, then $H(0)=0$, and every $h$ at which the
  instance is overloaded is mapped to 0, which is mapped to itself: all-miss attracts
  the overloaded hit rates.
\item If moreover $H(a)\ge a$ for some $a>0$, the least equilibrium is 0 and the
  greatest is at least $a$: the instance has two distinct extremal equilibria.
\end{enumerate}
\end{proposition}

\begin{proof}
(i) At $h=0$ every turn misses, the load is $\lambda\E[S^{\mathrm{miss}}]\ge1$, so
$W(0)=\infty$ and $H(0)=G(\infty;0)=0$. At an overloaded $h$, $H(h)=0$ in the same way,
and $H(0)=0$. (ii) 0 is an equilibrium, so $\underline h\le0$; and $a\in A$, so
$\bar h\ge a$ by \cref{L5:prop:tarski}.
\end{proof}

In a closed system the population bounds the wait and removes this route to $H(0)=0$,
but not the pool channel: if re-inserted contexts evict every paused prefix within a
think time, $H(0)=0$ still. Thrashing, a collapse of this kind cured by limiting the
load, is classical in paging systems; the admission cap of \cref{L5:sec:admit} is its
cure here.

\begin{example}[Two equilibria and a cap]\label{L5:ex:bistable}
Take deterministic $S^{\mathrm{hit}}=0.5$\,s, $S^{\mathrm{miss}}=5$\,s, $Z=7$\,s and the
wait channel alone, $G(T)=e^{-T/60\,\mathrm{s}}$, with the Pollaczek--Khinchine wait. At
$\lambda=0.25$/s the all-miss load is $1.25$, $H=0$ on $[0,2/9)$, and a grid search finds
three equilibria: 0 (stable), about 0.250 (unstable, slope $H'\approx12$) and about
0.882 (stable, slope 0.07). At $\lambda=0.18$/s the all-miss load is $0.9$ and the only
equilibrium is about 0.885. Lowering the load removes the collapse without moving the
good equilibrium much.
\end{example}

\begin{proposition}[The forced-miss multiplier]\label{L5:prop:multiplier}
Let $H(h)=c+kh$ near an equilibrium, with $c\ge0$ and $0\le k<1$, so $h_0=c/(1-k)$, and
force a fraction $\delta\in[0,1]$ of turns to miss, so the realised hit rate is
$(1-\delta)H(h)$. The new equilibrium is $h_\delta=(1-\delta)c/(1-(1-\delta)k)$, and
\[
  h_0-h_\delta=\frac{\delta\,c}{(1-k)(1-(1-\delta)k)}\ \ge\ \delta\,h_0 .
\]
\end{proposition}

\begin{proof}
$h_\delta(1-(1-\delta)k)=(1-\delta)c$ is the equation $h=(1-\delta)(c+kh)$. Over the
common denominator $(1-k)(1-(1-\delta)k)$ the numerator of $h_0-h_\delta$ is
$c(1-(1-\delta)k)-(1-\delta)(1-k)c=\delta c$. The inequality is
$1/(1-(1-\delta)k)\ge1$, true since $(1-\delta)k\ge0$.
\end{proof}

To first order in $\delta$ a forced miss converts a would-be hit with probability $h_0$
and induces $h_0k/(1-k)$ further misses on other turns, each with its own price, so its
marginal price is $h_0[\Phi_{\mathrm{forced}}+\tfrac{k}{1-k}\bar\Phi_{\mathrm{induced}}]$.
This cache multiplier multiplies the queue factor $(1-\rho_P)/(1-\rho_P')$ of
\cref{thm:price}; it does not replace it. A multiplier common to all sessions scales
every $u^P_i$ by the same factor: it leaves the eviction order unchanged and moves only
the comparisons with $\theta_P$ (keep or drop, admit).

\section{What is exact and what is approximate}\label{L5:sec:honest}

\begin{itemize}
\item \emph{Exact, given the model:} the price bracket (\cref{thm:price}), the prices at
  processor-sharing stations (\cref{prop:decode-price,L5:prop:link-price}), the
  monotonicity of the decode batch (\cref{thm:L-mono}), the finite-population
  inequalities (\cref{prop:finite}), the price-blind bounds (\cref{prop:blind}), and
  \cref{L5:prop:shadow,L5:prop:hold,L5:prop:tarski,L5:prop:collapse,L5:prop:multiplier}.
\item \emph{Modelling choices:} the Poisson reading of turn arrivals, which the session
  loop makes non-Poisson (\cref{prop:finite} shows the open formula overstates the wait
  for exponential work); equal slot sizes (with unequal contexts a turn needs several
  slots and the memory station has no closed form); the link as fair bandwidth sharing
  (a transfer engine that sends in arrival order would add a head-of-line term to
  \cref{L5:prop:link-price}); a single survival law and the mean wait in the feedback
  map; and the pricing of one eviction, a transient event, by a stationary comparison.
\item \emph{Estimates:} $p_i$, $\tau_i$ and $G$ must be predicted from a session's phase,
  tool and history; $a$, $b$, $x_0$, $\kappa$, $B$ and $\varphi$ must be measured on the
  instances.
\end{itemize}

\section{Exercises}

\begin{exercise}\label{L5:ex:split}
In \cref{L5:lem:split} let $\sigma_{\mathrm{hit}}=\sigma_{\mathrm{miss}}=\sigma$. Find the
hit rate that maximises the mixture term, and show that the mixture term exceeds the
append term exactly when $h(1-h)(\mu_{\mathrm{miss}}-\mu_{\mathrm{hit}})^2>\sigma^2$.
\emph{Hint:} $h(1-h)$ is maximised at $h=1/2$.
\end{exercise}

\begin{exercise}\label{L5:ex:compare}
A session with context $K$ is paused at both instances. Using \cref{L5:prop:shadow},
show that for $K$ large enough it is kept at the prefill instance whenever it is kept at
the decode instance, if $\theta_P=\theta_D$ and $b>0$.
\emph{Hint:} $\Phi_i\ge\Delta S_i\ge bK^2/2$ grows quadratically, $\Phi^X_i$ linearly in $K$.
\end{exercise}
